% ══════════════════════════════════════════════════════════════════════════════
%  EXPERIMENT 3 — BOUNCE EFFICIENCY                    (Lesson 8, T3 Week 3)
%  From Lesson 8 of the A4 Forces Lesson Booklet; answer keys 8.* kept.
% ══════════════════════════════════════════════════════════════════════════════
\experimentdiv{Bounce Efficiency}{Drop height and bounce height \textbullet{} Repeat trials and anomalies \textbullet{} Calculating efficiency}{8}\label{sec:3}
\lessoncodes{VC2S8U15 \textbullet{} VC2S8I01, I02, I03, I04, I05, I06}

\vocabhead
\begin{ctab}
\begin{tabular}{|SW{4.2cm}|W{11.6cm}|}
  \hline\hd & \thead{Word} & \thead{Definition (copy from the slide)} \\ \hline
  & \scaf{\textbf{Elastic potential energy}} & \scaf{The energy stored in an object when it is squashed or stretched.} \\ \hline
  & \textbf{Thermal energy} & \cell{8.v2} \\ \hline
  & \textbf{Drop height} & \cell{8.v3} \\ \hline
  & \textbf{Bounce height} & \cell{8.v4} \\ \hline
  & \textbf{Wasted energy} & \cell{8.v5} \\ \hline
  & \textbf{Anomaly} & \cell{8.v6} \\ \hline
\end{tabular}
\end{ctab}

\begin{lipanel}
\inquiryli{Measure heights precisely, spot anomalies and calculate with your results.}
\understandli{Use energy transformations to explain why a ball never bounces back to its drop height.}
\sctitle
\begin{itemize}[leftmargin=1.4em, itemsep=0pt, topsep=1pt]
  \item I can read the bounce height to the bottom of the ball, at eye level.
  \item I can leave an anomaly out of a mean and say why.
  \item I can calculate the efficiency of a bounce.
\end{itemize}
\end{lipanel}

\begin{skillsfocus}
\focus{QP1}{Write a hypothesis in the faculty frame.}
\focus{PC4}{Read each bounce at eye level, three times.}
\focus{PA1}{Record trials, means and efficiencies in a THUD table.}
\focus{PA2}{Plot the means and draw a line of best fit.}
\focus{PA3}{Calculate means and efficiencies.}
\focus{PA7}{Spot an anomaly and leave it out of the mean.}
\focus{EV3}{Name an error and how to reduce it.}
\end{skillsfocus}

% ─── READING ─────────────────────────────────────────────────────────────────
\exppart{Before the lab}{Teal}

\begin{readingbox}{Reading --- The Rules of the Bounce}
Every ball used in serious sport has to pass a bounce test. A tennis ball, for example, is dropped from a height of 254 centimetres onto concrete. To be approved for play, it must bounce back up to between 135 and 147 centimetres --- a little more than half the height it fell from.

Why doesn't a ball bounce all the way back up? As it falls, its gravitational potential energy becomes kinetic energy. When it hits the ground the ball squashes, and for a split second its energy is stored as elastic potential energy, like a squashed spring. The ball then springs back into shape and pushes itself back up. But some of the energy is wasted: the rubber warms up a little (thermal energy) and you hear the bounce (sound energy).

Squash players know this well. A cold squash ball hardly bounces at all, so players hit it hard for a few minutes before a match. As the ball warms up, the air inside it pushes harder on the rubber, and the rubber itself becomes springier, so the ball bounces higher.

In this experiment you will run your own bounce test, and use the heights to work out how efficient each bounce is.
\end{readingbox}

\rcreset
\rcqk{8.c1}{From what height is a tennis ball dropped in the bounce test, and how high must it bounce to be approved?}{2.2cm}
\rcqk{8.c2}{Explain, using the words \emph{elastic potential energy}, what happens to the ball's energy when it hits the ground.}{2.4cm}
\rcqk{8.c3}{Why can a ball never bounce back up to the height it was dropped from? Name the two forms of wasted energy.}{2.4cm}
\rcqk{8.c4}{Why do squash players hit the ball hard for a few minutes before a match?}{2.2cm}

% ─── TECHNIQUE ────────────────────────────────────────────────────────────────
\exppart{Technique spotlight}{Blue}

\begin{technique}{Reading a height at eye level}{PC4}
The top of a bounce lasts only a moment. Crouch so your eyes are level with where the ball will stop --- looking from above or below makes the reading wrong (parallax). Read to the \emph{bottom} of the ball every time, or film the bounce in slow motion and pause it at the highest point. One person drops, one person reads.
\end{technique}

\begin{technique}{Spotting an anomaly}{PA7}
An \textbf{anomaly} is a result that does not fit the pattern of the others: 52, 55 and \textbf{41}\,cm. Circle it, suggest what went wrong (the ball was pushed, or read from above), and leave it out of the mean --- or repeat that trial.
\end{technique}

\begin{infobox}{Bounce efficiency}
\[ \text{Efficiency} = \frac{\text{bounce height}}{\text{drop height}} \times 100\% \]
At the drop height and at the top of the bounce the ball is still, so all its energy is gravitational potential energy, which depends on height. Comparing the two heights compares the energy.
\end{infobox}

\begin{workedex}{Calculating bounce efficiency}
A ball is dropped from 100\,cm and bounces back up to 62\,cm.

\smallskip
\textbf{Efficiency:}\quad $62 \div 100 \times 100\% = \mathbf{62\%}$ \qquad \textbf{Wasted:}\quad $100\% - 62\% = \mathbf{38\%}$, as thermal energy and sound.
\end{workedex}

\begin{wedo}{Calculating bounce efficiency}
A tennis ball is dropped from 80\,cm and bounces back up to 44\,cm.
\pt{a}Find the efficiency of the bounce.
\wbox{8.wd1a}{2cm}
\pt{b}What percentage of the energy is wasted, and where does it go?
\wbox{8.wd1b}{2cm}
\end{wedo}

% ─── PLAN, SAFETY AND METHOD ─────────────────────────────────────────────────
\exppart{Plan, safety and method}{Teal}

\textbf{Aim:} To investigate the effect of changing the drop height of a ball on its bounce height.

\task Complete the variables table. \skr{PC2}
\begin{ctab}
\begin{tabular}{|SW{4.4cm}|W{10.2cm}|}
  \hline\hd & \thead{Variable} & \thead{In this experiment} \\ \hline
  & \scaf{Independent} & \scaf{The drop height: 40, 80, 120 and 160\,cm} \\ \hline
  & Dependent & \cell{8.t1-r2} \\ \hline
  & Controlled & \cell{8.t1-r3} \\ \hline
  & Controlled & \cell{8.t1-r4} \\ \hline
\end{tabular}
\end{ctab}

\needspace{8\baselineskip}
\task Write your hypothesis. Will the efficiency change as the drop height increases? \skr{QP1}

\scaf{If the drop height is increased, the bounce height will \dots}
\abox{8.pred}{2.5cm}

\subsubsection*{Equipment}
Tennis ball (or another rubber ball) \textbullet{} two metre rulers or a tape measure, fixed upright to a wall \textbullet{} sticky tape \textbullet{} hard, flat floor \textbullet{} phone with slow-motion video (optional)

\begin{center}
\diagBounceSetup
\expfigcap{Measure the drop height and the bounce height to the bottom of the ball, against rulers fixed to a wall.}
\end{center}

\needspace{12\baselineskip}
\task Complete the risk assessment. \skr{PC3}
\begin{ctab}
\begin{tabular}{|SW{4.0cm}|W{5.2cm}|W{5.4cm}|}
  \hline\hd & \thead{Risk (what could go wrong)} & \thead{Prevention} & \thead{Responding (if it happens)} \\ \hline
  & \scaf{Falling while reaching for a high drop} & \scaf{Keep every drop height within reach from the floor; never stand on chairs or desks} & \scaf{Tell the teacher straight away; first aid if needed} \\ \hline
  & Tripping over bags while chasing the ball & & \\ \hline
\end{tabular}
\end{ctab}

\subsubsection*{Method}
\begin{enumerate}[itemsep=2pt]
  \item Fix the metre rulers upright against the wall, with 0\,cm touching the floor.
  \item Hold the ball beside the rulers so that the \textbf{bottom} of the ball is level with the 40\,cm mark.
  \item Let go --- do not push or throw the ball.
  \item Read the height of the bottom of the ball at the top of its first bounce, at eye level. Record it in Table 3.1.
  \item Repeat steps 2--4 so that you have three trials.
  \item Repeat for drop heights of 80\,cm, 120\,cm and 160\,cm.
\end{enumerate}

% ─── RESULTS ─────────────────────────────────────────────────────────────────
\exppart{Results}{Amber}

\subsubsection*{Table 3.1: Bounce height and efficiency for each drop height \skr{PA1}}
\begin{ctab}
\begin{tabular}{|SZ{2.3cm}|Z{2.3cm}|Z{2.3cm}|Z{2.3cm}|Z{2.6cm}|Z{2.4cm}|}
  \hline\hd & \thead{Drop height (cm)} & \thead{Bounce, trial 1 (cm)} & \thead{Bounce, trial 2 (cm)} & \thead{Bounce, trial 3 (cm)} & \thead{Mean bounce (cm)} & \thead{Efficiency (\%)} \\ \hline
  & \scaf{e.g.\ 40} & \scaf{23} & \scaf{21} & \scaf{22} & \scaf{22} & \scaf{55} \\ \hline
  & 40  & \cell{8.res-r2c2} & \cell{8.res-r2c3} & \cell{8.res-r2c4} & \cell{8.res-r2c5} & \cell{8.res-r2c6} \\ \hline
  & 80  & \cell{8.res-r3c2} & \cell{8.res-r3c3} & \cell{8.res-r3c4} & \cell{8.res-r3c5} & \cell{8.res-r3c6} \\ \hline
  & 120 & \cell{8.res-r4c2} & \cell{8.res-r4c3} & \cell{8.res-r4c4} & \cell{8.res-r4c5} & \cell{8.res-r4c6} \\ \hline
  & 160 & \cell{8.res-r5c2} & \cell{8.res-r5c3} & \cell{8.res-r5c4} & \cell{8.res-r5c5} & \cell{8.res-r5c6} \\ \hline
\end{tabular}
\end{ctab}

\needspace{20\baselineskip}
\task Draw a line graph of your results. Put \textbf{drop height (cm)} on the horizontal axis and \textbf{mean bounce height (cm)} on the vertical axis. Plot each point, draw a line of best fit, and give your graph a title. \skr{PA2}
\graphgrid{9cm}

% ─── ANALYSIS ────────────────────────────────────────────────────────────────
\exppart{Analysis and evaluation}{Coral}

\task What happened to the bounce height as the drop height increased? Use two means from your table. \skr{PA5}
\abox{8.t2}{3cm}

\needspace{10\baselineskip}
\task Look at your efficiency column. Did the efficiency stay about the same, rise or fall as the drop height increased? Use numbers from your results.
\abox{8.t3}{3cm}

\needspace{8\baselineskip}
\task A ball dropped from 150\,cm bounces to 90\,cm. Calculate the efficiency of the bounce. \skr{PA3}
\wbox{8.t4}{2.6cm}

\needspace{8\baselineskip}
\task Another group's bounce heights for an 80\,cm drop were 52, 55 and 41\,cm. Which value is the anomaly? Find the mean of the other two. \skr{PA7}
\wbox{8.t5}{2.6cm}

\needspace{10\baselineskip}
\task Explain why the ball never bounces back up to its drop height. Use the words \emph{gravitational potential}, \emph{kinetic}, \emph{elastic potential} and \emph{thermal}.
\abox{8.t6}{4cm}

\needspace{10\baselineskip}
\task Identify one error in your method, and suggest how to reduce it. \skr{EV3}
\abox{8.t8}{3.5cm}

\needspace{12\baselineskip}
\task Write a conclusion: state the trend, back it with numbers from your table, and say whether your hypothesis was supported.

\scaf{As the drop height increased from 40\,cm to 160\,cm, the mean bounce height increased from \dots\,cm to \dots\,cm, while the efficiency stayed near \dots\,\%. This supports / does not support the hypothesis.}
\abox{8.concl}{5cm}

\begin{yr8box}
\begin{keepext}
\exttask Is your line of best fit straight, and does it pass through (or close to) the origin? Calculate its gradient, and explain how the gradient is connected to the efficiency of the ball.
\wbox{8.x8}{4cm}
\end{keepext}
\begin{keepext}
\exttask Design a fair test to find out whether the \emph{temperature} of a ball changes its bounce efficiency.
\pt{a}Name the independent and dependent variables.
\abox{8.x11a}{2cm}
\pt{b}Name two variables you would keep the same.
\abox{8.x11b}{2cm}
\pt{c}Describe a safe way to make the ball warmer and colder.
\abox{8.x11c}{2.2cm}
\end{keepext}
\end{yr8box}

% ─── SKILL CHECK ─────────────────────────────────────────────────────────────
\exppart{Skill check}{NavyBlue}

\begin{skillcheck}
\scrow{QP1}{My hypothesis names the drop height and the bounce height.}
\scrow{PC4}{I read each bounce at eye level, to the bottom of the ball.}
\scrow{PA1}{My table has a title, headings, units in the headings and matching decimals.}
\scrow{PA2}{I drew a line of best fit, not dot to dot.}
\scrow{PA3}{I calculated the means and efficiencies correctly.}
\scrow{PA7}{I spotted an anomaly and left it out of the mean.}
\scrow{EV3}{I named an error and a way to reduce it.}
\end{skillcheck}

\task Which skill will you focus on in Experiment 4? Explain what you will do differently.
\abox{8.sc}{5cm}
